\chapter{Metodología y diseño experimental}
\label{cap:metodologia}
\section{Diseño, unidad experimental y preguntas contrastables}
Se realizó un experimento factorial controlado sobre una oferta ficticia de analista de datos. La unidad profesional es el perfil base; sus seis CVs son observaciones emparejadas, no seis personas independientes. La variable objetivo $y^*$ expresa cumplimiento de una rúbrica conocida. La etiqueta observada $\tilde y$ representa la decisión histórica simulada que recibe el modelo. Esta separación permite detectar que reproducir fielmente una etiqueta histórica puede empeorar la evaluación profesional.

Se cruzaron cuatro escenarios principales, seis tratamientos y dos clasificadores, y se añadió un control del tamaño muestral para ambos modelos sin mitigación. El diseño se repitió sobre cinco corpus generados independientemente. Se mantuvieron comunes, dentro de cada réplica, los perfiles, las particiones y los números aleatorios que determinan retención y corrupción. Así, las comparaciones entre tratamientos son pareadas y no mezclan una intervención con un cambio arbitrario del conjunto de prueba.

Las hipótesis de trabajo son: la corrupción de etiquetas produce disparidades mayores que un desequilibrio aislado de representación; ocultar únicamente los datos explícitos no garantiza invariancia si quedan señales indirectas; la eliminación completa de las señales insertadas reduce las diferencias entre variantes, aunque puede conservar errores compartidos; y reparar parte de las etiquetas puede mejorar la recuperación de positivos. Estas expectativas guían los contrastes, pero no constituyen resultados. El protocolo quedó fijado antes del ajuste completo; no se realizó una preregistración externa.

@@DIAGRAMA_PROCESO@@

\section{Generación de perfiles y diversidad de los CVs}
\subsection{Variables profesionales y criterio de referencia}
Cada réplica contiene 2.000 perfiles, obtenidos mediante muestreo aleatorio y aceptación hasta reunir 1.000 positivos y 1.000 negativos. La probabilidad inicial de cada una de las cuatro competencias es 0,55; el equilibrio de clases se impone después mediante la rúbrica. En consecuencia, las frecuencias finales de competencias no tienen por qué mantener esa probabilidad ni ser independientes. No se pretende imitar la distribución real de solicitantes.

Se generan conocimientos de SQL, Python, visualización y estadística; experiencia relevante entre uno y ocho años; e inglés B1, B2 o C1. La puntuación es
\begin{equation}
 \begin{aligned}q&=2x_{\mathrm{SQL}}+2x_{\mathrm{Python}}+x_{\mathrm{visualización}}+x_{\mathrm{estadística}}+\mathbb{1}(e\geq2)+\mathbb{1}(i\in\{B2,C1\}),\\ y^*&=\mathbb{1}(q\geq6).\end{aligned}
\end{equation}
La mayor ponderación de SQL y Python representa una decisión de diseño de la oferta, no una escala validada por especialistas. La puntuación no se incluye en el texto que reciben los clasificadores. Las competencias ausentes no aparecen como negaciones; por ello, interpretar correctamente negaciones explícitas queda fuera del alcance del banco de pruebas. Una regla aplicada a los atributos estructurados alcanzaría acierto perfecto por definición: sirve para verificar etiquetas, no como competidor independiente de la tarea de lectura del CV.

Para evitar que el corpus se reduzca a seis textos casi idénticos, se incorporan seis titulaciones, ocho sectores, doce competencias complementarias, entre uno y tres equipos de trabajo, cinco duraciones de formación, seis proyectos, cinco entregables, tres modalidades y tres disponibilidades. Cada competencia puntuable dispone de tres formulaciones y se altera su orden. Se utilizan ocho disposiciones de secciones durante el desarrollo y se varía la longitud mediante una sección no puntuable opcional. Estas variables enriquecen el texto, pero no añaden puntos a la rúbrica. Son variaciones combinatorias controladas, no una evaluación experta del realismo de las trayectorias.

\subsection{Variantes demográficas y señales indirectas}
Por cada perfil se generan seis variantes: hombre y mujer, cruzados con tres intervalos de edad, 24--34, 38--48 y 52--62 años. Se muestrea una edad en cada intervalo y se comparte entre las variantes de género. Se seleccionan nombres ficticios de listas de 16 nombres por categoría y doce apellidos. Su eventual coincidencia con personas existentes es accidental; no se consultó información personal.

Todas las variantes mantienen competencias, experiencia, titulación, proyecto, estilo y longitud estructural. Cambian el nombre, el género y la edad de la cabecera, la flexión \emph{organizado/organizada} y el año de finalización, calculado como $2026-\mathrm{edad}+20$. El año introduce deliberadamente una señal de edad fuera de la cabecera. La finalización uniforme a los veinte años simplifica trayectorias educativas diversas, y el CV no detalla una cronología completa de empleos. La experiencia es experiencia relevante acumulada, no antigüedad total; no se infiere ausencia de actividad laboral durante los años restantes. Estas simplificaciones limitan la plausibilidad biográfica.

El emparejamiento representa intervenciones sobre \emph{señales textuales} demográficas. No permite afirmar que se haya construido un modelo causal completo del efecto de género o edad sobre una carrera. El género binario y tres intervalos separados de edad son decisiones instrumentales; se excluyen identidades no binarias y edades fuera de esos rangos. La equidad observada solo se refiere a las categorías representadas.

@@DIAGRAMA_CV@@

\subsection{Prueba de cambio de formato}
Además del texto principal, se crea una versión de prueba con cuatro órdenes no utilizados en desarrollo, nuevos títulos de secciones y dos sustituciones de formulación. Se mantiene la cualificación y la etiqueta de cada perfil. Esta condición, denominada formato no visto, mide sensibilidad a cambios moderados de presentación. No incorpora un nuevo idioma, documentos escaneados, errores de extracción ni un vocabulario profesional completamente diferente. Por tanto, se evita denominarla validación externa.

\section{Escenarios de entrenamiento y referencia de evaluación}
La tabla~\ref{tab:probabilidades} fija las probabilidades del mecanismo de sesgo. La retención se decide con un uniforme $u_r$ por CV: se incluye si $u_r<r_g$. La corrupción utiliza otro uniforme independiente $u_c$ y cambia únicamente positivos de $y^*=1$ a $\tilde y=0$ cuando $u_c<c_g$. No se convierten negativos en positivos. Las probabilidades son parámetros artificiales de estrés; no proceden de estadísticas de discriminación.

\begin{table}[htbp]\centering\small
\caption{Parámetros de sesgo por intersección. H: hombre; M: mujer.}\label{tab:probabilidades}
\begin{tabular}{lrrrrrr}\toprule
Grupo & H 24--34&H 38--48&H 52--62&M 24--34&M 38--48&M 52--62\\\midrule
Retención $r_g$&1,00&0,90&0,65&0,70&0,50&0,25\\
Corrupción $c_g$&0,05&0,15&0,40&0,20&0,35&0,60\\\bottomrule
\end{tabular}\end{table}

El escenario control conserva todas las observaciones y etiquetas. Representación aplica solo retención; etiquetado aplica solo corrupción; combinado aplica ambas. El quinto escenario, control de tamaño, toma tantos CVs como el combinado, distribuidos casi uniformemente entre doce celdas de grupo y clase verdadera, con una diferencia máxima de una observación. Sus etiquetas son correctas. Este control ayuda a separar cantidad total y composición, pero no iguala exactamente la cobertura de perfiles profesionales: no identifica de forma perfecta un efecto causal puro de representación.

La validación y la prueba permanecen completas, equilibradas y etiquetadas con $y^*$. Esta referencia limpia es una ventaja deliberada del estudio sintético. Una organización con etiquetas históricas dudosas no dispone automáticamente de ella; necesitaría una muestra revisada y un criterio profesional defendible. Seleccionar hiperparámetros contra etiquetas limpias puede producir estimaciones más favorables que seleccionarlos contra decisiones históricas sesgadas.

\section{Separación de datos y ajuste de los modelos}
Se utilizan semillas 11, 22, 33, 44 y 55; la semilla de generación es 20261013 más la semilla de réplica. En cada corpus, el 60\% de perfiles se asigna a entrenamiento, el 20\% a validación y el 20\% a prueba, estratificando por $y^*$. Antes de retener CVs hay 7.200 documentos de entrenamiento, 2.400 de validación y 2.400 de prueba. Todas las variantes del mismo perfil permanecen en la misma partición; también su versión de formato alternativo. No se utiliza la prueba para escoger C, el tratamiento o el umbral.

Los textos se representan mediante TF--IDF de unigramas y bigramas, frecuencia mínima de dos documentos, máximo de 14.000 características, frecuencia sublineal y normalización L2. Se normalizan acentos y mayúsculas, conservando números y tokens de un carácter. El vocabulario y los pesos IDF se aprenden exclusivamente en entrenamiento y se reutilizan en validación y prueba. Cada representación textual y escenario ajusta su propio vectorizador. Los métodos que comparten representación y documentos reutilizan esa matriz, evitando diferencias accidentales.

Se entrenan regresión logística (LR), con \texttt{liblinear}, y SVM lineal (SVM), con formulación primal. Ambos utilizan regularización L2, hasta 10.000 iteraciones y tolerancia $10^{-5}$ \cite{sklearnapi}. Se prueban $C\in\{0{,}1,1,10\}$, seleccionando mayor F1 de validación; ante empate se conserva el menor C. Los avisos de convergencia en los ajustes principales se tratan como errores. El umbral principal es margen cero, equivalente a probabilidad 0,5 en LR; el margen SVM no se interpreta como probabilidad calibrada.

Estas familias permiten inspeccionar pesos y aportaciones léxicas con un coste moderado. No se entrenan redes neuronales, modelos generativos ni representaciones preentrenadas. El trabajo estudia mecanismos de aprendizaje supervisado textual, no el comportamiento de todos los sistemas comerciales de reclutamiento.

\section{Mitigación por componentes y auditoría parcial}
Se comparan seis tratamientos. El original utiliza todo el texto y $\tilde y$. La reponderación mantiene texto y etiquetas, pero asigna
\begin{equation}
 w(g,y)=\frac{n_g n_y}{n\,n_{g,y}},
\end{equation}
donde las frecuencias se calculan solo en entrenamiento y respecto a la etiqueta utilizada por ese tratamiento. Los pesos tienen media uno. Se exige que ninguna celda esté vacía. Este procedimiento reduce asociación entre grupo y etiqueta \cite{kamiran2012}, pero no identifica qué observaciones están mal etiquetadas ni restaura automáticamente el tamaño de los grupos.

La supresión de cabecera elimina la primera línea completa. La limpieza completa añade reglas textuales que neutralizan la flexión de género y retiran el año de finalización en ambos formatos. No consulta etiquetas ni metadatos para transformar un CV. Sin embargo, conoce todas las señales insertadas por el generador: constituye un caso favorable de limpieza exhaustiva y no un anonimizador general de CVs reales.

El quinto tratamiento combina limpieza y reponderación. El sexto añade reparación de etiquetas: se selecciona al azar el 30\% de los perfiles de entrenamiento, 360 por réplica, sin examinar errores ni pertenencia demográfica para elegirlos. Para las variantes retenidas de esos perfiles, se sustituye $\tilde y$ por $y^*$. Después se recalculan los pesos y se entrena con el texto limpio. El coste expresado en perfiles no equivale necesariamente al 30\% exacto de documentos retenidos. Esta intervención simula acceso a una auditoría perfecta en la muestra elegida; no se llevó a cabo una revisión humana ni se midieron desacuerdos entre evaluadores.

@@DIAGRAMA_MITIGACION@@

La comparación por componentes permite distinguir qué aporta retirar identificadores, neutralizar señales indirectas, modificar la distribución de pesos y corregir información supervisora. Si la limpieza hace idénticas las seis entradas de un perfil, un clasificador determinista necesariamente produce seis decisiones idénticas. Esta propiedad se verifica, pero no se interpreta como prueba de justicia general. Todos los candidatos de un perfil podrían recibir una decisión errónea común.

\section{Métricas, incertidumbre y pruebas de sensibilidad}
Se evalúan F1, precisión, sensibilidad, exactitud, área ROC y precisión media. La métrica principal de disparidad es $\Delta_{TPR}=\max_gTPR_g-\min_gTPR_g$ sobre los seis grupos. Se acompañan brecha de FPR, diferencia máxima de tasas de selección y el máximo entre brechas de TPR y FPR. La comparación de grupos se hace contra $y^*$, nunca contra las etiquetas corrompidas.

Las pruebas emparejadas calculan la proporción de decisiones distintas al cambiar género con la misma edad (tres pares por perfil), y al cambiar intervalo de edad manteniendo género (seis pares). Sus denominadores son pares, no candidatos independientes. AUC describe ordenación de puntuaciones; no prueba ausencia de discriminación en un umbral operativo. La F1 depende de la prevalencia y no mide directamente coste económico o daño social.

Se presentan medias y desviaciones estándar de cinco corpus independientes. Las diferencias entre tratamientos se calculan primero dentro de cada réplica. Los intervalos percentiles del 95\% remuestrean con reemplazo las cinco diferencias completas, 10.000 veces. Son intervalos bootstrap exploratorios con pocas unidades; no se atribuye precisión poblacional a miles de CVs correlacionados ni se declara significación tras múltiples comparaciones. La inferencia queda condicionada al generador y a las cinco realizaciones disponibles.

Tres pruebas complementarias examinan robustez. Primero, se evalúan los formatos reservados con los modelos ya ajustados. Segundo, LR con $C=1$ se reentrena con multiplicadores de corrupción 0, 0,5, 1 y 1,5, manteniendo retención y números aleatorios. Tercero, se reentrenan limpieza con pesos y auditoría con pesos bajo umbrales de rúbrica 5 y 7 y bajo seis requisitos con peso igual y umbral 4. En este último análisis se conserva la partición original, sin reestratificar por la rúbrica alternativa. No son nuevos resultados del mismo clasificador: cada condición requiere un ajuste.

Como diagnóstico adicional, se exploran 61 umbrales entre $-3$ y $3$ en validación para LR original y con auditoría. Se elige el de mayor F1, favoreciendo el más próximo a cero en empate, y se evalúa una sola vez en prueba. No se buscan umbrales específicos por grupo. El protocolo, las versiones, las salidas individuales y los controles de integridad se conservan en el material reproducible citado en la bibliografía \cite{repositorio2026}.
